\subsection{字典序乘积}

设 \((\mathbf{E}_t)_{i\in I}\) 是一族有序集，由良序集 \(I\) 索引。考虑乘积集 \(\mathbf{E} = \prod_{i\in I}\mathbf{E}_t\) ，且关系

`` \(x \in E\) 并且 \(y \in E\) ，并且对于最小索引 \(t \in I\) 使得 \(\mathrm{pr}_t x \neq \mathrm{pr}_t y\) ，我们有 \(\mathrm{pr}_t x < \mathrm{pr}_t y\) '',

，我们将其记为 \(\mathbf{R}\{x, y\}\) 。显然有 \(\mathbf{R}\{x, x\}\) 等价于 \(x \in \mathbf{E}\) ，且 \(\mathbf{R}\{x, y\}\) 蕴含 \(\mathbf{R}\{x, x\}\) 并且 \(\mathbf{R}\{y, y\}\) ，以及 \((\mathbf{R}\{x, y\}\) 并且 \(\mathbf{R}\{y, x\})\) 蕴含 \(x = y\) 。此外容易验证 \((\mathbf{R}\{x, y\}\) 并且 \(\mathbf{R}\{y, z\})\) 蕴含 \(\mathbf{R}\{x, z\}\) （考虑最小索引 \(t \in I\) ，使得三个元素 \(\mathrm{pr}_t x\) , \(\mathrm{pr}_t y\) , \(\mathrm{pr}_t z\) 中至少有两个不相等）；因此 \(\mathbf{R}\{x, y\}\) 是乘积集合E上的一个序关系。该关系及其所定义的序被称为E上的字典序关系和字典序（由 I 上及 \(\mathbf{E}_t\) 上的给定序所诱导）；具有该序的集合 E 称为有序集族 \((\mathbf{E}_t)_{t \in I}\) 的字典序乘积。如果每个 \(\mathbf{E}_t\) 是全序的，则字典序乘积也是全序的。

